% Discussion (1,000 words target)

\subsection{Constitutional Feasibility}
% Weighted voting question:
% Does fractional vote-casting satisfy "one person, one vote"?
% Each person's vote still counts equally (county aggregate weight proportional)
% Precedent: Shareholder voting (proportional to shares)
%
% Legal analysis:
% Likely requires Supreme Court case (no precedent for federal weighted voting)
% May require constitutional amendment (Article I, Section 2 interpretation)

\subsection{Practical Implementation}
% House procedural changes:
% Electronic voting systems (fractional vote counting)
% Committee assignments (fractional members?)
% Quorum calculations (fractional thresholds)
%
% Feasibility: Technically straightforward, politically challenging

\subsection{Advantages Over Algorithmic Redistricting}
% 1. Eliminates redistricting: No decadal boundary disputes
% 2. Respects existing communities: Counties are established political units
% 3. Transparent: No algorithm complexity
% 4. Dynamic: Automatically adjusts to population shifts (re-weight annually?)

\subsection{Disadvantages}
% 1. Weighted voting novelty: No federal precedent
% 2. County inequality: Some counties poorly designed (historical artifacts)
% 3. No VRA compliance: Counties may dilute minority voting strength
% 4. Prisoner populations: Counties with prisons get inflated populations (unless adjusted)

\subsection{Future Work}
% Test with alternative units (municipalities, media markets)
% Analyze VRA compliance implications
% Survey public attitudes (weighted voting acceptability)
